% JETSPIN parallelization documentation

\section{Parallelization\label{sec:Parallelization}}

The parallel infrastructure of JETSPIN incorporates the necessary
data distribution and communication structures. The parallel strategy
underlying JETSPIN is the Replicated Data (RD) scheme\cite{smith1991molecular},
where fundamental data of the simulated system are reproduced on all
processing nodes. In simulations of nanofibers, the fundamental data
consist of position, velocity, and viscoelastic stress arrays of each
bead in which the nanojet is discretised (see below Cap \ref{chap:JETSPIN-Model-and}).
Further data defining mass and charge of each bead are also replicated.
However, all auxiliary data are distributed in equal portion of data
(as much as possible) for each processor. Despite being available
other parallel strategies such as the Domain Decomposition\cite{brown1993domain},
our experience has shown that such volume of data is by no means prohibitive
on current parallel computers.

By the RD scheme, we adopt in JETSPIN a simple communication strategy
of the communication between nodes, which is handled by global summation
routines. The module \texttt{version\_mod} located in the \texttt{parallel\_version\_mod.f90}
file contains all the global communication routines which exploit
the MPI (Message Passing Interface) library. Note that a FORTRAN90
compiler and an MPI implementation for the specific machine architecture
are required in order to compile JETSPIN in parallel mode. An alternative
version of the module \texttt{version\_mod} is located in the \texttt{serial\_version\_mod.f90}
file, and it can be easily selected by appropriate flags in the \texttt{makefile}
at compile time. By selecting this version, JETSPIN can also be run
on serial computers without modification, even though the code has
been designed to run on parallel computers.

The size system is strictly time-dependent as mentioned in Sec \ref{sec:Structure}
and, therefore, the memory of various service bookkeeping arrays is
dynamically distributed over all the processing nodes. In particular,
the bookkeeping array size, declared as \texttt{mxchunk}, is managed
by the routine \texttt{set\_mxchunk}. All the beads of the discretised
nanofiber are assigned at every time step to a specific node by the
routine \texttt{set\_chunk}, and their temporary data are stored in
bookkeeping arrays belonging to the assigned node. It is worth stressing
that the communication latency makes the parallelization efficiency
strictly dependent on the system size. Therefore, we only advice the
use of JETSPIN in parallel mode whenever the user expects a system
size with at least 50 beads for each node (further details in Ref
\cite{lauricella2015jetspin}).
